\chapter*{ABSTRACT}
\begin{center}
  \large
  \textbf{Evaluation of \textit{Multi-Object Tracking }Algorithms and Video Analytics Platforms for Measuring Consumer Behavior in Commercial Areas}
\end{center}
% Menyembunyikan nomor halaman
\thispagestyle{empty}

\begin{flushleft}
  \setlength{\tabcolsep}{0pt}
  \bfseries
  \begin{tabular}{lc@{\hspace{6pt}}l}
  Student Name / NRP&: &Danendra Galang Yugastama / 5024231049\\
  Department&: &Computer Engineering FTEIC - ITS\\
  Advisor&: &1. Ir.Arta Kusuma Hernanda, S.T., M.T.\\
  & & 2. Wernher von Braun, S.T., M.T.\\
  \end{tabular}
  \vspace{4ex}
\end{flushleft}
\textbf{Abstract}

Video analytics is increasingly crucial for the retail industry to automatically extract consumer behavior metrics such as people counting, dwell time, and queue time. However, visual occlusion from visitor density triggers identity switches, which distorts the accuracy of business metrics. Since high academic tracking metrics do not necessarily guarantee low practical errors, this research proposes a comparative evaluation of four contemporary Multi-Object Tracking (MOT) algorithms: ByteTrack, BoT-SORT, OC-SORT, and Deep OC-SORT. The evaluation utilizes a Tracking-by-Detection architecture with a YOLO11 detector on real-world operational CCTV footage with high occlusion levels. The performance of each algorithm is measured by comparing the system's output against manual observation data (ground truth) and a commercial platform. This research aims to identify the most robust tracking algorithm in minimizing Mean Absolute Error (MAE) and Mean Absolute Percentage Error (MAPE) for retail analytics. Furthermore, this study analyzes the statistical correlation between academic metrics (MOTA and IDF1) and real-world precision to prove the relevance of theoretical benchmarks to the practical needs of the retail industry.

\vspace{2ex}
\noindent
\textbf{Keywords: \emph{Identity Switches, Multi-Object Tracking, Occlusion, Retail Analytics, YOLO11.}}